\documentclass[../../main.tex]{subfiles}
\begin{document}


\begin{problem}
对作业排序问题证明：当且仅当子集合\( J \)中的作业可以按下述规则处理时，\(J\)才表示一个可行解。

~\\
即如果 J 中的作业还没有分配处理时间，则将它分配在时间片\([\alpha -1,\alpha]\)处理，其中\(\alpha\)是使得\(1 \leq r \leq d_i\)的最大整数 \(r\) ，且时间片\([\alpha -1,\alpha]\)是空的。
\end{problem}
\begin{solution}

\textbf{必要性（如果 \( J \) 是可行解，那么可以按规则处理）}

如果\(J\)是一个可行解,那么\(J\)中的作业是可以按照截止时间 \( d_i \) 的非降次序排列的。

对于其中的任意作业\(J_k\),如果其时间期限为\(d_k\)且时间片\([d_k-1,d_k]\)为空,则可以分配.
若非空,则向前寻找一个空时间片\([\alpha-1,\alpha]\),其中\(\alpha\)是使得\(1 \leq r \leq d_i\)的最大整数 \(r\)。

由于\(J\)是可行解,那么一定可以找到.

~\\
\textbf{充分性（如果按规则处理所有作业，则 \( J \) 是可行解）}

当所有作业都按照规则分配处理时间时，每个作业的处理时间不会超过其截止时间 \( d_i \)。
易得\(J\)是一个可行解。


综上,命题得证。
\end{solution}
\end{document}